\chapter{Results}

$\hat{q}(t)$ and $\diff E / \diff x (t)$ were evaluated following Section  \ref{sec:numerical_computation} using different parameter values of $Q_{s1/2}$, $m$ and $\mathbf{v}$. Unless otherwise stated, the same value was used for both saturation scales $Q_{s1} = Q_{s2}$, denoted simply by $Q_s$. For the remaining parameters, the same value was always used: $g(\mu)=1$, $N_c = 3$, $N_f = 2$. The magnitude of the velocity, $v$, was kept at $v=1$. All quantities were calculated either in the $0 < t\leq \qty{0.1}{fm}$ time interval or the $0 < t \leq \qty{0.6}{fm}$ time interval.

Before presenting those results, however, the energy density and temperature calculated within the calculation of $\diff E / \diff x$ are shown.

\section{Energy density and ``Temperature''}
The event-averaged energy density as a function of proper time, $\left< \varepsilon (\tau) \right>$, given by Eq. (\ref{eq:Pablos_energy_density}), was evaluated for $0 < \tau < \qty{0.1}{fm}$, for several values of the saturation scales, as illustrated in Fig. \ref{fig:energy_Q_sweep}. $\left< \varepsilon (\tau) \right>$ was also evaluated for a fixed $Q_s$, varying $m$ so as to obtain different values of the $Q_s/m$ ratio, as illustrated in Fig. \ref{fig:energy_Qm_sweep}. As suggested by Eqs. (\ref{eq:Pablos_energy_density}) and (\ref{eq:regularized_MV_G_h}), $\left< \varepsilon (\tau) \right>$ scales with $Q_s^4$ and the speed of its decay increases with the ratio $Q_s/m$.

The value of $Q_s$ which reproduces \cite{Carrington_2022}'s initial energy density, $\left<\varepsilon_0 \right>$, is $Q_s = \qty{3.8079}{GeV}$. They used a saturation constant of $Q_s = \qty{2.0}{GeV}$, and the difference is justified by the fact that another regularization method of the gluon densities was used in their case. As such, when comparing these results with \cite{Carrington_2022}'s, the $Q_s = \qty{3.8079}{GeV}$ curves should be used.

The corresponding temperatures, $\left<T(\tau) \right>$, given by Eq. (\ref{eq:final_full_T}), are illustrated in Fig. \ref{fig:temperature_Q_sweep}. As $\left<T \right>$ is proportional to the fourth root of the energy density, the $\left< \varepsilon \right>$ range of $\sim [4,20](\qty{}{GeV^4})$ produced a narrower range of $\left<T \right>$, $\sim [0.9,1.2](\qty{}{GeV})$. As mentioned in Section \ref{sec:temperature}, the temperature should be much greater than the mass of the flavors of quarks included in $N_f$. This is the case here, as the masses of the up- and down-quarks are less than $\qty{10}{MeV}$. The strange-quark, with a mass of the order of $\qty{100}{MeV}$, is excluded from the \gls{QGP}. \cite{PDG_2024}

\begin{figure}
\centering
\begin{tikzpicture}
\begin{axis}[
	cycle list={
        [colors of colormap={875,750,625,500,375,250,125,0} of hot]},
	%cycle list={[of colormap=hot]},
	%cycle list name = resultstyle,
	title = {$\left< \varepsilon (\tau) \right> \quad (m = \qty{0.2}{GeV})$},
	scaled ticks=false,
    xticklabel style={
        /pgf/number format/fixed,
        /pgf/number format/fixed zerofill,
        /pgf/number format/precision=2
        },
    xlabel = $\tau \; (\qty{}{fm})$,
    ylabel = $\left<\varepsilon \right> \; (\qty{}{GeV^4})$,
    %	ylabel = {\shortstack{$\left<\varepsilon_0 \right>$\\$(\qty{}{GeV^2})$}},
   	%ylabel style={rotate=-90},
   	xmin = 0, xmax = 0.1, ymin = 0,
   	%ytick distance = 0.5,
   	no markers,
   	 xticklabel={
        \ifdim\tick pt=0pt
        \else
            \pgfmathprintnumber{\tick}
        \fi
    },
    yticklabel={
        \ifdim\tick pt=0pt
        \else
            \pgfmathprintnumber{\tick}
        \fi
    },
    legend pos=outer north east
    ]
    \addlegendimage{empty legend}
	\addlegendentry{\hspace{-0.7cm} $Q_s \; (\qty{}{GeV})$}
\foreach \Q in {4.0000,3.8079,3.6000,3.4000,3.2000,3.0000,2.8000,2.6000} {
    \edef\filename{Data/energy_sweeping_Q/MV_Energy_Q\Q.txt}
    \addplot table [x index=0, y index=1] {\filename};
	\addlegendentryexpanded{$\pgfmathprintnumber[fixed,fixed zerofill,precision=2]{\Q}$}	
	%\pgfmathprintnumberto[fixed,fixed zerofill,precision=2]{\Q}{\roundedQ}
	%\edef\temp{\noexpand\addlegendentry{$\roundedQ$}}
	%\temp
}
	\addplot[gray, densely dotted] {15.9773};

\end{axis}
\end{tikzpicture}
\caption{$\left< \varepsilon (\tau) \right>$ for several values of $Q_s$. Dotted in gray, \cite{Carrington_2022}'s value for the initial energy density, $\left<\varepsilon_0\right> = \qty{15.9773}{GeV^4}$; it matches this work's initial energy density for $Q_s \approx \qty{3.8079}{GeV}$.}
\label{fig:energy_Q_sweep}
\end{figure}

\begin{figure}
\centering
\begin{tikzpicture}
\begin{axis}[
	title = {$\left< \varepsilon (\tau) \right> \quad (Q_s = \qty{3.8079}{GeV})$},
	scaled ticks=false,
    xticklabel style={
        /pgf/number format/fixed,
        /pgf/number format/fixed zerofill,
        /pgf/number format/precision=2
        },
    xlabel = $\tau \; (\qty{}{fm})$,
    ylabel = $\left<\varepsilon \right> \; (\qty{}{GeV^4})$,
    %	ylabel = {\shortstack{$\left<\varepsilon_0 \right>$\\$(\qty{}{GeV^2})$}},
   	%ylabel style={rotate=-90},
   	xmin = 0, xmax = 0.1,
   	ymin=0,
   	%ytick distance = 0.5,
   	no markers,
   	 xticklabel={
        \ifdim\tick pt=0pt
        \else
            \pgfmathprintnumber{\tick}
        \fi
    },
    yticklabel={
        \ifdim\tick pt=0pt
        \else
            \pgfmathprintnumber{\tick}
        \fi
    },
    legend pos=outer north east
    ]
    \addlegendimage{empty legend}
	\addlegendentry{\hspace{-0.7cm} $Q_s/m$}
\foreach \Q in {5.0000,10.0000,20.0000,40.0000} {
    \edef\filename{Data/energy_sweeping_Qm/MV_Energy_Q\Q.txt}
    \addplot table [x index=0, y index=1] {\filename};
	\addlegendentryexpanded{$\pgfmathprintnumber[fixed,fixed zerofill,precision=0]{\Q}$}	
	%\pgfmathprintnumberto[fixed,fixed zerofill,precision=2]{\Q}{\roundedQ}
	%\edef\temp{\noexpand\addlegendentry{$\roundedQ$}}
	%\temp
}
	\addplot[gray, densely dotted] {15.9773};

\end{axis}
\end{tikzpicture}
\caption{$\left< \varepsilon (\tau) \right>$ for several values of $Q_s/m$.}
\label{fig:energy_Qm_sweep}
\end{figure}

\begin{figure}
\centering
\begin{tikzpicture}
\begin{axis}[
	cycle list={
        [colors of colormap={875,750,625,500,375,250,125,0} of hot]},
	%cycle list name = resultstyle,
	title = {$\left< T (\tau) \right> \quad (m = \qty{0.2}{GeV})$},
	scaled ticks=false,
    xticklabel style={
        /pgf/number format/fixed,
        /pgf/number format/fixed zerofill,
        /pgf/number format/precision=2
        },
    xlabel = $\tau \; (\qty{}{fm})$,
    ylabel = $\left< T \right> \; (\qty{}{GeV})$,
    %	ylabel = {\shortstack{$\left<\varepsilon_0 \right>$\\$(\qty{}{GeV^2})$}},
   	%ylabel style={rotate=-90},
   	xmin = 0, xmax = 0.1,
   	%ytick distance = 0.5,
   	no markers,
   	 xticklabel={
        \ifdim\tick pt=0pt
        \else
            \pgfmathprintnumber{\tick}
        \fi
    },
    yticklabel={
        \ifdim\tick pt=0pt
        \else
            \pgfmathprintnumber{\tick}
        \fi
    },
    legend pos=outer north east,
    clip=false
    ]
    \addlegendimage{empty legend}
	\addlegendentry{\hspace{-0.7cm} $Q_s \; (\qty{}{GeV})$}
\foreach \Q in {4.0000,3.8079,3.6000,3.4000,3.2000,3.0000,2.8000,2.6000} {

    \edef\filename{Data/temperature_sweeping_Q/MV_Temperature_Q\Q.txt}

    \addplot table [x index=0, y index=1] {\filename};
	\addlegendentryexpanded{$\pgfmathprintnumber[fixed,fixed zerofill,precision=2]{\Q}$}
}
\node[anchor=west,scale = 0.8]
    at (axis description cs:1,0.67)
    {$Q_s (\qty{}{GeV})$};
\end{axis}
\end{tikzpicture}
\caption{$\left< T (\tau) \right>$ for several values of $Q_s$. }
\label{fig:temperature_Q_sweep}
\end{figure}

\subsection{Dependence on $Q_s$}
Figures \ref{fig:qHat_Q_sweep} and \ref{fig:dEdx_Q_sweep} illustrate $\hat{q}$ and $\diff E / \diff x$ as functions of time for different values of $Q_s$. Overall, the values of $\hat{q}$ at $t = \qty{0.1}{fm}$ are of the order $\qty{10}{GeV^2/fm}$ and those of $\diff E / \diff x$ are of the order $10^{-1} \qty{}{GeV/fm}$ \textcolor{red}{(expected values)}. $\hat{q}$ plateaus at around $t = \qty{0.5}{fm}$.

As predicted by \cite{Carrington_2022}'s expansions of $\hat{q}$ and $\diff E / \diff x$ in powers of proper time, when $v_\parallel = 0$, the initial slope of the stopping power is $0$, unlike that of $\hat{q}$. Furthermore, the initial slope of the $Q_s = \qty{3.8079}{GeV}$ curve agrees with \cite{Carrington_2022}'s result. 

\begin{figure}
\centering
\input{Graphs/qHat_sweeping_Q}
\caption{$\hat{q}(t)$ for several values of $Q_s$. \textcolor{red}{Since $v_\parallel=0$, $t=\tau$.} }
\label{fig:qHat_Q_sweep}
\end{figure}

\begin{figure}
\centering
\begin{tikzpicture}
\begin{axis}[
	cycle list={
        [colors of colormap={875,750,625,500,375,250,125,0} of hot]},
	%cycle list name = resultstyle,
	title = {$\diff E / \diff x \quad (v_\perp=1, m = \qty{0.2}{GeV})$},
	scaled ticks=false,
    xticklabel style={
        /pgf/number format/fixed,
        /pgf/number format/fixed zerofill,
        /pgf/number format/precision=2
        },
    xlabel = $t \; (\qty{}{fm})$,
    ylabel = $-\diff E/ \diff x \; (\qty{}{GeV/fm})$,
    %	ylabel = {\shortstack{$\left<\varepsilon_0 \right>$\\$(\qty{}{GeV^2})$}},
   	%ylabel style={rotate=-90},
   	xmin = 0, xmax = 0.1, ymin = 0,
   	%ytick distance = 0.5,
   	no markers,
   	 xticklabel={
        \ifdim\tick pt=0pt
        \else
            \pgfmathprintnumber{\tick}
        \fi
    },
    yticklabel={
        \ifdim\tick pt=0pt
        \else
            \pgfmathprintnumber{\tick}
        \fi
    },
    legend pos=outer north east
    ]
    \addlegendimage{empty legend}
	\addlegendentry{\hspace{-0.7cm} $Q_s \; (\qty{}{GeV})$}
\foreach \Q in {4.0000,3.8079,3.6000,3.4000,3.2000,3.0000,2.8000,2.6000} {
    \edef\filename{Data/dEdx_sweeping_Q/MV_dEdx_Q\Q.txt}
    \addplot table [x index=0, y index=1] {\filename};
	\addlegendentryexpanded{$\pgfmathprintnumber[fixed,fixed zerofill,precision=2]{\Q}$}	
	%\pgfmathprintnumberto[fixed,fixed zerofill,precision=2]{\Q}{\roundedQ}
	%\edef\temp{\noexpand\addlegendentry{$\roundedQ$}}
	%\temp
}

\end{axis}
\end{tikzpicture}
\caption{$\diff E / \diff x$ for several values of $Q_s$. \textcolor{red}{Since $v_\parallel=0$, $t=\tau$.} }
\label{fig:dEdx_Q_sweep}
\end{figure}

\begin{figure}
\centering
\input{Graphs/qHat_0.6}
\caption{$\hat{q}$ up to $t = \qty{0.6}{fm}$. \textcolor{red}{Since $v_\parallel=0$, $t=\tau$.} }
\label{fig:qHat_0.6}
\end{figure}

